\section*{Author integer-complexity and integrality context}
\paragraph{Definition 11 (Straight Line Program).} An SLP (straight line program) $P$ (for computing an integer) of length (or size) $n$ is a sequence of integers $a_0,\ldots,a_n$ with $a_0=1$ and $a_k=a_i\circ a_j$ with $i,j<k$ for $\circ\in\{+,-,\times\}$. We say that the SLP $P$ computes the integer $a_n$. For an integer $N$, we define the straight line complexity $\tau(N)$ of $N$ to be the length of the smallest SLP computing $N$.
\paragraph{Constant-free circuit complexity.} In the same spirit, one can define constant-free circuit complexity of polynomials where the given constants belong to the set $\{-1,0,1\}$.\footnote{To use $2^n$ in the circuit, one has to build up a circuit for $2^n$, of size $\log n$, from 1; whereas in the usual sense of circuit size, constants are free. Thus, $f_d:=2^{2^d}x^d$ has $O(\log d)$-size circuit but requires $\Omega(d)$-size constant-free circuit.} We denote, $\tau(f)$ as the size of the minimal constant-free circuit computing $f$. Trivially, $L(f)\le\tau(f)$.
\paragraph{Definition 28 (Easy sequence).} A sequence $(a_n)_n$ of integers is said to be ``easy to compute'' if there exists a polynomial $p$ such that straight line complexity of $a_n$, i.e. $\tau(a_n)\le p(\log n)$, for $n\ge1$. If a sequence is not easy to compute, it is said to be hard. In fact, for most numbers $N$, one can show that $\tau(N)\ge\log N/\log\log N$ (``close'' to the trivial upper bound) [14,35]. It is believed that $(d!)$ is hard to compute. In fact, its hardness is deeply connected to the infamous integer factoring problem. [42] showed that $d!$ being easy to compute will imply factoring is easy in the non-uniform setting.\footnote{However, this result does not imply that natural numbers can be factored in polynomial time in the Turing-Machine model, as the numbers used can be $\mathrm{poly}(n)$-bits.}
\paragraph{Definition 33 (Ultimately easy).} A sequence of integers $(a_n)$ is ultimately easy if there exists another sequence $(b_n)$ such that $\tau(a_nb_n)\le\mathrm{poly}(\log n)$ for all large enough $n$.
\paragraph{Definition 34 (Ultimate complexity).} Define the ultimate complexity of an integer $n$ as the minimum $\tau$-complexity of its multiple, i.e. $\tau_1(n)=\min_{b\in\mathbb Z\setminus\{0\}}\tau(b\cdot n)$. It is clear that Definition 33 can be stated wrt $\tau_1$. We remark that $\tau_1(n_1n_2)\le\tau_1(n_1)+\tau_1(n_2)+1$, for any $n_1,n_2\in\mathbb Z$.
\paragraph{Binomial series.} For rational $n$, $(x+a)^n=\sum_{k\ge0}\binom nkx^ka^{n-k}$, where $\binom nk=n(n-1)\cdots(n-k+1)/k!$. A formal power series is $A=\sum_{i\ge0}A_ix^i$; its degree-$d$ truncation is $\mathrm{trunc}(A,d)=\sum_{0\le i\le d}A_ix^i$.
\paragraph{Fact 74 (Folklore).} Product of any $k$ consecutive positive integers is divisible by $k!$.
\paragraph{Definition 75 ($p$-adic valuation).} Let $p$ be a prime and $n\in\mathbb Z$. We denote $p$-adic valuation of $n$ as $\nu_p(n)$ to be the highest exponent such that $p^{\nu_p(n)}\mid n$. Formally, $\nu_p:\mathbb Z\longrightarrow\mathbb N$ defined by $\nu_p(n)=\max\{v\in\mathbb N:p^v\mid n\}$. Note that, by definition, $\nu_p(\mathrm{rad}(n))=1$ if $p\mid n$, and 0 otherwise.
\paragraph{Theorem 76 (Legendre's formula).} For a prime $p$ and $n\in\mathbb N$, $\nu_p(n)=\sum_{j=1}^\infty\lfloor n/p^j\rfloor$.
\paragraph{Theorem 77 (Restatement of Theorem 31, Integral power series).} Let $k\in\mathbb N$. Define $f_k(x):=(1+k^2x)^{1/k}$. Then, $f_k(x)\in\mathbb Z[[x]]$.
\begin{proof}By binomial expansion, $f_k\in\mathbb Q[[x]]$. Let $f_k(x)=\sum_{d\ge0}a_dx^d$. We'll prove by strong induction that indeed the coefficients are integers. Obviously $a_0=1$, and assume that for $m\in\mathbb N$ we have proved that $a_\ell\in\mathbb Z$ for $0\le\ell<m$. The coefficient at $x^m$ in $(\sum_{d=0}^\infty a_dx^d)^k=(1+\sum_{d=1}^\infty a_dx^d)^k$ is equal to $k\cdot a_m$ plus a bunch of terms that we know are integer by the induction hypothesis; hence $k\cdot a_m=b\in\mathbb Z$. But by the binomial series formula we have
\[a_m=k^{2m}\binom{1/k}m=\frac{k^{2m}\prod_{j=0}^{m-1}(1/k-j)}{m!}=\frac{k^m\prod_{j=0}^{m-1}(1-kj)}{m!}.\]
It suffices to prove that $k\mid b$. If we can show that $\nu_p(b)\ge\nu_p(k)$ for every prime $p$ dividing $k$, this would certainly imply that $k\mid b$. So, fix a prime $p\mid k$. Note that
\[b=k\cdot a_m=X/m!,\qquad X:=k^{m+1}\prod_{j=0}^{m-1}(1-kj).\]
As, $p\mid k$, we must have $\prod_{j=0}^{m-1}(1-kj)\equiv1\bmod p$. Thus, $\nu_p(X)=\nu_p(k^{m+1})=(m+1)\nu_p(k)$. And by Theorem 76, $\nu_p(m!)=\sum_{j=1}^\infty\lfloor m/p^j\rfloor<\sum_{j=1}^\infty m/p^j=m/p-1\le m$. Thus,
\[\nu_p(b)=\nu_p(X)-\nu_p(m!)\ge(m+1)\nu_p(k)-m\ge\nu_p(k),\]
as we wanted. Putting it together gives $a_m\in\mathbb Z$ proving the inductive step. Hence, the conclusion follows.\end{proof}
It suffices to prove the integrality of $(1+k^2x)^{1/k}$ for the candidate functions $(1+k^2x)^{i/k}$ with $i<k$, as in the authors' Appendix G. In Section6 the authors work throughout with $(1+k^2x)^{i/k}$, $i,k\in\mathbb N$ and $i<k$; Theorem31 states their integrality.
